\subsection{Discrete-time Markov Decision Process}

In this subsection, I explore the Discrete-time Markov Decision Process (MDP) as a mathematical framework for modeling sequential decision-making problems within the context of RL. I focus on how the foundational elements of MDPs can be applied to economic decision-making problems. The subsequent discussion will detail the structure and components of MDPs, highlighting their relevance and adaptation to the economics domain.

RL uses an agent-environment interface as the basis for its mathematical notation. The \enquote{agent} refers to the entity responsible for selecting the control variables or actions, paralleling the role of economic agents in economic theory. Conversely, \enquote{environment} is used about the system that simulates forward the choice variables and returns the next state and reward; in the context of this thesis, this role is fulfilled by the IAM. Figure (\ref{fig:agent-environment interaction}) shows how the dynamics of agent-environment interactions in RL are typically modeled. Please see \textcite[Chapter 3]{Sutton2018intro_to_rl} for a thorough introduction. 

\begin{figure}[]
    \centering
    \includesvg[width=12cm]{images/ch3/agent-env-interaction.svg}
    \caption{Conventional modeling of agent-environment interaction.}
    \label{fig:agent-environment interaction}
\end{figure}

A Markov Decision Process (MDP) is a way to model sequential decision-making. Using the same notation as Sutton and Barto \textcite[Chapter 3]{Sutton2018intro_to_rl}, the MDP is defined by the tuple \(\mathcal{M} = \{\mathcal{S}, \mathcal{A}, \{P_{s, a}\}, \gamma, R\}\). Here, \(\mathcal{S}= \{s_1, s_2, \dots, s_n\}\) represents a finite set of \(n\) states with the current state at any time \(t\) expressed as the vector \(\bm{s}_t \in \mathbb{R}^n\). The set \(\mathcal{A} = \{a_1, a_2, \dots, a_k\}\) is a finite set of \(k\) actions with the action at time \(t\) represented by the vector \(\bm{a}_t \in \mathbb{R}^k\). The function \( P_{s, a}(\cdot) \) denotes the probability of transitioning to a subsequent state \(s'\), upon taking action \(a\) in state \(s\). \(\gamma \in (0,1)\) denotes the discount factor. Lastly, the reward function \(R: \mathcal{S} \times \mathcal{A} \rightarrow  \mathbb{R} \) assigns a real-valued reward to each state-action pair.

This framework offers a generalized approach to modeling sequential decision-making in economics. To illustrate this connection more explicitly, consider this simple intertemporal optimization problem in economics. Suppose the state is represented by capital \(s_t = K_t\) and that the economic agent at each time step \(t\) chooses the consumption level for the time step, meaning that the action is \(a_t= C_t\). The agent is forward-looking and has time preferences to discount future utility \(u(C_t) = r_t(s_t, a_t)\) with the discount factor \(\gamma = \beta\). Additionally, the agent must consider uncertainty through transition probabilities in stochastic economic models.

In RL, the agent chooses actions according to a policy. For the purposes of this introduction, I will use a simplified notation, where the policy function is defined as a direct mapping from the set of states \(\mathcal{S}\) to the set of actions \(\mathcal{A}\) represented as \(\pi: \mathcal{S} \rightarrow \mathcal{A}\)\footnote{This represents a simplification from the typical stochastic policy functions in reinforcement learning (which you will find throughout \textcite{Sutton2018intro_to_rl} as well), where a policy is normally defined as a mapping from states and actions to a probability distribution: \(\pi: \mathcal{S} \times \mathcal{A} \rightarrow [0,1]\). That means \(\pi(a_t|s_t) = Pr(A_t=a_t | S_t=s_t)\).}. The value function for a policy \(\pi\) can be evaluated at any state \(s_t\) and is given by:

\begin{equation}
    V_{\pi}(s) = \mathbb{E}_\pi\left[ \sum_{k=0}^{\infty} \gamma^k R_{t+k+1} \mid S_t = s \right], \quad \forall s \in \mathcal{S},
\end{equation}

where the expectation \(\mathbb{E}_\pi\) is over state transitions and rewards, and the reward function is expressed as \(R_t\) for brevity. This expression evaluates the sum of expected discounted rewards obtained by following the policy \(\pi\) from state \(s\). This recursive relationship is foundational to the Bellman equation for the value function. Introduced in the seminal work in \textcite{BELLMAN1957AMDP}, the Bellman equation can be expressed as follows:

\begin{equation}
    V_\pi(s_t) = \mathbb{E}_\pi [R_{t+1} + \gamma V(s_{t+1})]
\end{equation}

The objective of RL is to identify a policy that maximizes the rewards across all states, maximizing the value function. This is encapsulated in the principle of Bellman optimality, which can be mathematically represented as:

\begin{equation}
    \pi^*(s) = \argmax_{a} \mathbb{E} [R_{t+1} + \gamma V_\pi(S_{t+1}) \mid S_t = s, A_t = a],
\end{equation}

where \(\pi^*\) denotes the optimal policy. To simplify and make my notation more compact, I will also introduce the action-value function \(Q_\pi(\bm{s}_t, \bm{a}_t)\), which is the right-hand side of the Bellman equation:

\begin{equation}
    Q_\pi(s, a) = \mathbb{E} [R_{t+1} + \gamma V_\pi(S_{t+1}) \mid S_t = s, A_t = a]
\end{equation}

This function quantifies the expected return starting from state \(s\), taking action \(a\), and thereafter following policy \(\pi\). 